\chapter*{Conclusion Générale}
\markboth{Conclusion Générale}{}
\addcontentsline{toc}{chapter}{Conclusion Générale}

L'introduction posait la question de savoir si un même système pouvait tenir ensemble la plomberie opérationnelle de la gestion des actifs industriels, utilisateurs, machines, ordres de travail, interventions, et une couche prédictive à part entière, sans que l'une ne tire l'autre vers le bas. C'est le cas : les sept modules ML (P1 à P7) s'insèrent dans le même cycle de vie des ordres de travail et d'intervention qu'un technicien utilise déjà, plutôt que d'y être greffés, et chacune de leurs sorties s'accompagne d'une explication en langage clair plutôt que d'un score de modèle brut, ce qui était la condition pour que la plateforme inspire assez confiance pour être réellement utilisée.

Y parvenir a suivi un chemin volontairement peu spectaculaire : d'abord les besoins, puis une architecture modulaire, puis une réalisation par tranches, puis des tests confrontés à ce que les besoins exigeaient réellement, plutôt que de se précipiter directement sur les modèles. C'est cet ordre qui a empêché la couche ML de devenir une démonstration greffée sur un système EAM fonctionnel, plutôt qu'une partie intégrante de celui-ci.

Sur le plan technique, ce stage a imposé une confrontation réelle avec les frontières des microservices, avec l'application du ML à des problématiques qui ne ressemblent pas à un jeu de données de manuel (régression fondée sur la survie, détection d'anomalies, réconciliation de plusieurs sources d'évidence en désaccord), et avec la construction d'un pipeline de génération augmentée par récupération (RAG) de bout en bout plutôt que le simple câblage d'une solution existante. Sur un plan moins technique, cela a signifié porter seul un périmètre de six mois : décider de la prochaine étape, en documenter les raisons, et refuser de considérer un composant comme terminé tant qu'il n'avait pas été réellement confronté à son comportement réel plutôt que supposé fonctionner.

Ce qui reste à faire en dit autant que ce qui a été livré. Le service RAG devrait accepter davantage de formats de documents qu'aujourd'hui. Les modèles sont entraînés sur des données simulées et ont besoin de véritables résultats confirmés par des techniciens avant que leurs chiffres d'exactitude n'aient un sens, une limite énoncée clairement plutôt que dissimulée au Chapitre~5. Et l'automatisation des pièces détachées devrait rester exactement aussi contrôlée par l'humain qu'elle l'est aujourd'hui, sinon davantage, à mesure qu'elle mûrit.
